% BHU PYQ -- Transcribed & Corrected
% Paper: MATMJ-21 / MATMN-21: Algebra (Mid-Term)
% Exam: B.A. / B.Sc. (Hons.) Semester II Examination 2024-25 (NEP)
% Department: Department of Mathematics

\documentclass[11pt,a4paper]{article}
\usepackage[a4paper,top=2.5cm,bottom=2.5cm,left=2.8cm,right=2.8cm,headheight=14pt]{geometry}
\usepackage{amsmath,amssymb}
\usepackage[T1]{fontenc}
\usepackage[utf8]{inputenc}
\usepackage{lmodern,microtype,fancyhdr,enumitem,hyperref,lastpage,parskip}

\hypersetup{colorlinks=true,urlcolor=black,linkcolor=black,pdftitle={MATMJ21/MATMN21: Algebra Mid-Term -- BHU B.Sc. Sem II 2024-25 (NEP)}}
\pagestyle{fancy}\fancyhf{}
\fancyhead[L]{\small   \textit{Banaras Hindu University}}
\fancyhead[R]{\small   \textit{MATMJ-21 / MATMN-21: Algebra (Mid-Term)}}
\fancyfoot[C]{\small Page  \thepage\ of \pageref{LastPage}}


\newlist{parts}{enumerate}{2}
\setlist[parts,1]{label=(\alph*),leftmargin=*,itemsep=4pt,topsep=2pt}

\newcommand{\pts}[1]{\hfill   \textit{[#1]}}


\begin{document}


\begin{center}
  {\large\bfseries BANARAS HINDU UNIVERSITY}\\[2pt]
  {\normalsize B.A. / B.Sc. (Hons.) Semester II Mid-Term Test, 2024-25 (NEP)}\\[6pt]
  \rule{0.8\linewidth}{0.5pt}\\[6pt]
  {\large\bfseries DEPARTMENT OF MATHEMATICS}\\[4pt]
  {\large\bfseries Paper Code: MATMJ-21 / MATMN-21 --- Algebra}\\[4pt]
  {\normalsize Time: 1 Hour\hfill Maximum Marks: 20}
\end{center}

\rule{\linewidth}{0.4pt}\smallskip



\noindent   \textit{Instructions: Attempt any   \textbf{FOUR} questions. All questions carry equal marks (5 marks each).}
\bigskip




\noindent   \textbf{Question 1}

\begin{parts}
  \item Define $n$-th root of unity. Prove that if $\omega 
eq 1$ is an $n$-th root of unity, then $1 + \omega + \omega^2 + \dots + \omega^{n-1} = 0$.\pts{3}
  \item Let $z = 2 - 2i$. Find $|z|$, $ \text{Arg}(z)$, $ar{z}$, and $\sqrt{z}$. Plot $z$, $ar{z}$, and $\sqrt{z}$ on the complex plane.\pts{2}
\end{parts}

\medskip\rule{\linewidth}{0.2pt}\medskip




\noindent   \textbf{Question 2}

\begin{parts}
  \item Let $p(x)$ be a polynomial with real coefficients. Prove that if $\alpha = x + iy$ is a root of $p(x)$ with $y 
eq 0$, then $ar{\alpha} = x - iy$ is also a root of $p(x)$.\pts{3}
  \item Find the smallest degree polynomial $p(x)$ with real coefficients which has roots $3$ and $3 + 2i$.\pts{2}
\end{parts}

\medskip\rule{\linewidth}{0.2pt}\medskip




\noindent   \textbf{Question 3}

\begin{parts}
  \item Let $a$ and $b$ be two non-zero integers, and let $d = \gcd(a, b)$. Prove that there exist integers $x$ and $y$ such that $d = ax + by$.\pts{3}
  \item Using Euclidean algorithm, find $\gcd(8991, 12345)$.\pts{2}
\end{parts}

\medskip\rule{\linewidth}{0.2pt}\medskip




\noindent   \textbf{Question 4}

\begin{parts}
  \item Let $p(x) = ax^3 + bx^2 + cx + d$ be a cubic polynomial and let $\alpha, \beta, \gamma$ be its roots. Find the relation between the roots and coefficients of $p(x)$.\pts{3}
  \item Consider the polynomial $p(x) = x^3 - 5x^2 + 9x - 5$. If $\alpha, \beta, \gamma$ are the roots of $p(x)$, find the value of $\alpha + \beta + \gamma$, $\alphaeta + \beta\gamma + \gamma\alpha$, $\alphaeta\gamma$, and $\alpha^2 + \beta^2 + \gamma^2$.\pts{2}
\end{parts}

\medskip\rule{\linewidth}{0.2pt}\medskip




\noindent   \textbf{Question 5}

\begin{parts}
  \item Define equivalence relation.\pts{2}
  \item Let $m$ be a fixed positive integer. Define a relation $R$ on the set of integers $\mathbb{Z}$ such that $a R b \iff m \mid (a - b)$. Check whether $R$ is an equivalence relation or not. If yes, for $m = 5$, find the equivalence class of $2$.\pts{3}
\end{parts}

\vfill\rule{\linewidth}{0.4pt}

\begin{center}\small\href{https://bhu-exam.vercel.app/}{Click here to visit our website}\end{center}
\end{document}
